% Results I — full prose v1.

\subsection{Most measured breakthroughs are budget artifacts}

The restart control changes what the probe data mean. In the development
cohort, Ministral-3 exhibited apparently clean mid-trajectory crossings ---
until a budget-sensitivity re-probe showed the same problems solved from
16-token prefixes once the continuation budget was raised to 4{,}096
tokens: the ``breakthroughs'' marked where solutions began to fit the
budget, not where the prefix accumulated value. Applied to all 178 cells,
the frozen taxonomy yields the map in Figure~\ref{fig:regimemap}: for
Ministral-3, 33 cells are budget-limited, 17 instant, 8 terminal, 2
no-crossing, and 29 unsolved; for Gemma-3n, 26 instant, 4 budget-limited,
3 no-crossing, 55 unsolved --- and exactly \emph{one} prefix-limited cell
in the entire dataset (a development problem with $\Tv = 896$ and advantage
$0.75$). The 98 expansion cells, collected after the rules were frozen,
contributed zero new prefix-limited cases. Loosening the margin to
$\delta = 0.25$ adds four crossings, tightening to $0.75$ keeps two, and
the conservative $\Rhat$ envelope agrees with the primary labels on all
178 cells. Genuine prefix-locked value, at this scale and in these models,
is rare to absent.

\begin{figure}[t]\centering
\includegraphics[width=\linewidth]{a5_project_regime_map}
\caption{The final regime map: 89 problems $\times$ 2 models classified
under the frozen rules (top: per-cell strip grouped by cohort; bottom:
per-model totals). One cell of 178 is prefix-limited.}
\label{fig:regimemap}\end{figure}

\subsection{Restart dose--response separates failure modes}

The restart curves themselves carry the cleanest model-level finding
(Figure~\ref{fig:dose}). On expansion cells that are not instant, the
share of problems Ministral-3 solves from scratch climbs monotonically
with budget --- $0\% \to 19\% \to 45\% \to 79\%$ across $1{,}024 \to
8{,}192$ tokens, with zero non-monotone curves among 42 --- while
Gemma-3n stays flat at $10$--$12\%$ across the same budgets, with five
problems whose restart rate \emph{decreases} at larger budgets
(overthinking on natural problems, complementing constructed-task inverse
scaling \citep{gema2025inverse}). One model's failures are
\emph{compute-starved}; the other's are \emph{capability-limited}. The two
models agree on a collapsed regime class for only 42 of 89 problems:
small models do not share a single ``small-model'' failure mode, which
also cautions against averaging such models in evaluations.

\begin{figure}[t]\centering
\includegraphics[width=\linewidth]{wave3_restart_dose_response}
\caption{Restart dose--response on non-instant expansion cells.
Individual curves (left) and share solved from scratch (right):
Ministral-3 climbs $0\%\to79\%$; Gemma-3n stays flat.}
\label{fig:dose}\end{figure}

\subsection{Accumulated reasoning is compute compression}

The 13 threshold-censored Ministral-3 expansion cells --- trajectories
whose crossing was replicated only at the final anchor --- allow the
sharpest statement of what a long prefix is worth
(Figure~\ref{fig:matched}). At \emph{matched total compute}, continuing
the model's own prefix beats restarting in 11 of 13 cells (median
advantage $+0.31$); yet an $8{,}192$-token restart catches up in 11 of 13.
Accumulated reasoning therefore acts as \emph{compression} --- it reaches
the same solutions with fewer total tokens --- rather than as expanded
reachability. This instantiates, empirically and per problem, the no-benefit branch of
the dichotomy of \citet{wolf2026theory} --- absent self-reflection that
reliably localizes early errors, conditioning on past attempts offers no
asymptotic benefit over independent restarts (their analysis concerns
attempt-conditioned search; our truncated-prefix continuation is the
analogous object) --- and it
gives the sequential-versus-parallel debate \citep{snell2024,
ghosal2025mirage, sharma2025sequentialedge} a quantity it has lacked: the
value of a partial trace expressed in fresh-token units.

\begin{figure}[t]\centering
\includegraphics[width=.9\linewidth]{wave3_matched_compute}
\caption{Matched-compute value of long prefixes on 13 threshold-censored
cells: the prefix wins at equal total compute (11/13), and a large restart
budget closes the gap (11/13).}
\label{fig:matched}\end{figure}

\subsection{A genuinely stochastic solvability band}

Enlarging all 148 ambiguous cells to eight attempts resolved 55 as
effectively unsolvable ($\le 2/8$) and 38 as solvable ($\ge 6/8$) --- and
left \emph{55 genuinely intermediate} ($3$--$5$ of $8$;
Figure~\ref{fig:band}). At these anchors the continuation success
probability really is $\approx 0.3$--$0.6$: a property of the state, not
of the measurement. Intermediate states are therefore not binary
solvable/unsolvable, which undermines the single-crossing, monotone-value
picture implicit in binary-search process-labeling schemes
\citep{luo2024omegaprm} and in snowball-error accounts of monotone
degradation \citep{gan2025snowball}, and which converges with
trajectory-level evidence that recoverable and structural failures are
distinct populations \citep{islah2026}.

\begin{figure}[t]\centering
\includegraphics[width=\linewidth]{wave3_topup_resolution}
\caption{What ambiguous $1$--$3/4$ cells became at eight attempts: about a
third remain genuinely intermediate.}
\label{fig:band}\end{figure}
